\usepackage[german]{babel}
\usepackage[utf8]{inputenc}
\usepackage{graphpap,
		latexsym,
		color,
		amssymb,
		graphicx,
		amsfonts,
		enumerate,
		longtable,
		lscape,
		afterpage,
		}

\usepackage[hyphens,allowmove,obeyspaces
%,T1
]{url} % allows breaks behind hyphens

% must be here as it needs to be behind the other includes
\usepackage[a4paper,
		twoside=true,
		headheight=2cm,
		text={16.15cm,26cm}]{geometry}

\usepackage[pdftitle={Kommentiertes Vorlesungsverzeichnis},
% 		breaklinks=true,
% % 		bookmarks=true,
% % 	 	bookmarksnumbered, 
% % 		bookmarksopen=true,
% % 		bookmarksopenlevel=1,
% 		colorlinks=true,
% 		linkcolor=black,  % Farbe der Links (black,red,green,blue,...)
% 		a4paper,
% 		pdffitwindow=true,
% 		pdfpagelayout=SinglePage,
% 		dvips,
% 		pdftex, % so wird link aktiv in pdf-Datei
	   ]{hyperref}
\usepackage{breakurl} % hyperref muss vorher geladen sein

\setlength{\parindent}{0pt}


%% Definition der Strukturen %%
%
\newcommand{\chapter}[1]{{%
    \cleardoublepage%
    %~\vfill
    ~\\[40ex]
    {\Huge\bfseries\centerline{#1}}\par\nobreak
    %~\vfill
    \fakesection{#1}
       \newpage
}}
%
\newcommand{\btabt}{\large%
      \noindent\begin{tabular}{@{}p{3.8cm}@{~}p{12.1cm}@{}}
}
%
\newcommand{\etabt}{
     \end{tabular}
     \vspace{1mm}\hrule\vspace{5mm}
}
%
\newcommand{\btabb}{\vfill\vfill
  \vspace{1mm}\hrule\vspace{2mm}
  \normalsize\noindent
  \begin{tabular}{@{}p{5.5cm}@{~}p{10.4cm}@{}}
}
%
\newcommand{\etabb}{\end{tabular}}
%
\newcommand{\blit}{
    \vspace*{5mm}%
    \normalsize
    \textbf{Literatur:}\vspace*{-2mm}
    \begin{enumerate}[1.)]%
	\setlength{\parskip}{0pt}%
	\setlength{\itemsep}{0pt}
}
%
\newcommand{\elit}{\end{enumerate}}



\newcommand{\NN}{N.\,N.}
\newcommand{\nV}{n.\,V.}
\newcommand{\zB}{z.\,B.}

%% Definition von Professoren-Namen
%
\newcommand{\Ruzicka}{R{\r u}{\v z}i{\v c}ka}


%% einzelne Abteilungen
%

%% Definition Uni-Logo und Linie darunter
%
\newcommand{\UniLogo}{%
\raisebox{-1.0cm}{%
\fboxsep0pt%
\hyperlink{IHV}{%
\includegraphics[viewport= 0 0 145 175,height=2.7cm]{L-ALU-ab-2009-farbig}}}
\\[-1.25cm]
\rule{13.5cm}{2pt}}%
%
\newcommand{\UniLogoFront}{%
\raisebox{-1.8cm}{%
\fboxsep0pt{%
%\includegraphics[viewport= 0 0 2430 2880,height=4.7cm]{L-ALU}}}
\includegraphics[viewport= 0 0 243 288,height=4.7cm]{L-ALU-ab-2009-farbig}}}
\\[-1.6cm]
\rule{12.0cm}{3pt}}
%
\newcommand{\UniLogoMathFak}{%
\raisebox{-1.0cm}{%
\fboxsep0pt%
\hyperlink{IHV}{%
\includegraphics[viewport= 0 0 145 175,height=2.7cm]{L-ALU-ab-2009-farbig}}}
\\[-.95cm]
\rule{13.5cm}{2pt}}%


%% Definition der Koepfe %%
%
\newcommand{\newllogo}[5]{\expandafter\newcommand\csname#1\endcsname{{%
#2%
\begin{minipage}[b]{#3}%
\large\textbf{\textsf{#4}\mbox{}}
\end{minipage} \hfill %
\Semester\qquad#5\vspace{10mm}}}
}
%
\newcommand{\newlogo}[4]{\newllogo{#1}{#2}{#3}{#4}{\UniLogo}}
\newcommand{\newmlogo}[4]{\newllogo{#1}{#2}{#3}{#4}{\UniLogoMathFak}}

% Abteilung f"ur Angewandte Mathematik (\iam)
\newlogo{iam}{%
\includegraphics[height=1.5cm]{L-AM}\vspace{2ex}
\hspace*{1.3cm} % <- Dieses Leerzeichen ist notwendig...
}{5cm}{Abteilung f{\"u}r\\Angewandte Mathematik\\[-0.7em]}
%
% Abteilung f"ur Mathematische Logik (\ifl)
\newlogo{ifl}{%
\includegraphics[height=1.5cm]{L-Logik}\vspace{2ex}
\hspace*{2.3cm} % <- hier auch...
}{4.5cm}{Abteilung f{\"u}r\\ Mathematische Logik\\[-.7em]}
%
% Abteilung f"ur Mathematische Stochastik (\ifs)
\newlogo{ifs}{%
\includegraphics[height=1.5cm]{L-MS}\vspace{2ex}
\hfill
}{3cm}{}
%
% Abteilung f{\"u}r Reine Mathematik (\mai)
\newlogo{mai}{%
\rule{0pt}{1.5cm}\vspace{2ex}%
}{5cm}{%
Abteilung f{\"u}r\\ Reine Mathematik\\[-.7em]}
%
% Abteilung f{\"u}r Didaktik der Mathematik (\mad)
\newlogo{mad}{%
\rule{0pt}{1.5cm}\vspace{2ex}%
}{5cm}{%
Abteilung f{\"u}r\\Didaktik der Mathematik\\[-0.7em]}
%
% Didaktik -- PH Freiburg (madph)
\newlogo{madph}{%
\includegraphics[height=1.6cm]{Definition/L-PH}\hspace{3.0cm}\includegraphics[height=1.5cm]{Definition/L-PH-FACE}\vspace{2ex}%
}{1cm}{%
\\[-0.7em]}
%
% Abteilung f{\"u}r Finanzmathematik (\fin)
\newlogo{fin}{%
\rule{0pt}{1.5cm}\vspace{2ex}%
}{5cm}{%
Abteilung f{\"u}r\\Finanzmathematik\\[-0.7em]}
%
% Abteilung f{\"u}r Quantitative Finanzmarktforschung (\qff)
\newlogo{qff}{%
\rule{0pt}{1.5cm}%
\vspace{2ex}%
% \hspace*{4.0cm}%
}{7.5cm}{%
Abteilung f{\"u}r\\Quantitative Finanzmarktforschung\\[-0.7em]}
%
% Abteilung Medizinische Biometrie und Statistik am
% Institut f"ur Med. Biometrie und Med. Informatik (\imb)
\newlogo{imb}{%
\includegraphics[height=1.5cm]{L-imbi}\vspace{2ex}%
\hfill%
}{5.6cm}{%
Institut f{\"u}r\\Medizinische Biometrie und\\
Statistik\\[-0.7em]}
%
% Computer Vision Group (\cvg)
\newlogo{cvg}{%
\raisebox{0.8ex}{\includegraphics[height=.5cm]{Definition/L-Vision-2}}\vspace{2ex}
\hspace*{1.3cm} % <- Dieses Leerzeichen ist notwendig...
}{5cm}{Computer Vision Group\\[-0.7em]}
%
% Vorsitzender Pruefungsausschuesse (\pruef)
\newmlogo{pruef}{}{10cm}{%
Mathematisches Institut\\
Vorsitzender der 
Pr\"ufungsaussch\"usse Mathematik\\
Prof. Dr. S. Goette\\[-.7em]}
%
%

% Math. Inst. (\minst)
\newlogo{minst}{%
\rule{0pt}{1.5cm}\vspace{2ex}%
}{5cm}{%
Mathematisches Institut\\[-.7em]}
%

% Mathematisches Institut (\fak)
\newmlogo{fak}{}{5cm}{%
Mathematisches Institut\\[-.7em]}
%
% Mathematische und Physikalische Fakult"at (\mfak)
\newlogo{mfak}{%
\rule{0pt}{2.5cm}%
}{6.3cm}{%
\begin{center}%
Fakult\"at f\"ur\\%
Mathematik und Physik\\%
Mathematisches Institut\\%
\end{center}}
%
% Leere Kopfzeile mit Unilogo   (\leer)
\newmlogo{leer}{%
\rule{0pt}{1.5cm}%
}{10cm}{%
\phantom{Leer g}\\[-.5em]}


%% Definition der Sprechstundenseite %%
%
\newcommand{\Sprechstunden}{%
\fak \\ %
\bigskip\bigskip
{\renewcommand{\baselinestretch}{2.0}\tiny\large

\textbf{\Large Sprechstundenverzeichnis}\\[4ex]
Bitte informieren Sie sich \"uber die aktuell g\"ultigen Sprechstunden\\[2ex] der Professorinnen, der Professoren und Assistent(inn)en auf der Internet-Seite\\[4ex]
\url{http://home.mathematik.uni-freiburg.de/personen/list/sprechstunden.de.html}.}
}


%% Definition des Impressums %%
%
\newcommand{\impressum}{%
\thispagestyle{empty}%
\fakesection{Impressum}%
~\vfill
\hyperlink{IHV}{\textbf{Impressum}}\\[2ex]
Herausgeber: \\[1ex]
Mathematisches Institut \\
Eckerstr. 1\\
79104 Freiburg\\
Tel.: 0761-203-5534\\
E-Mail: institut@math.uni-freiburg.de
}
%





%
%
%% Definition der Lehrveranstaltungen %%
%
\newcommand{\vorlesung}[1]{%
	\kopfeintrag{Vorlesung}{#1}{#1}
}
\newcommand{\praktikum}[1]{%
	\kopfeintrag{Prakt. \"Ubung zu}{{#1}}{{#1}}
}
\newcommand{\praktischeuebung}[1]{%
	\kopfeintrag{Prakt. \"Ubung zu}{{#1}}{{#1}}
}
\newcommand{\proseminar}[1]{%
	\kopfeintrag{Proseminar}{#1}{#1}
}
\newcommand{\seminar}[1]{%
	\kopfeintrag{Seminar}{#1}{#1}
}
\newcommand{\forschungsseminar}[1]{%
	\kopfeintrag{Forschungsseminar}{#1}{#1}
}
\newcommand{\arbeitsgemeinschaft}[1]{%
	\kopfeintrag{Projektseminar}{#1}{#1}
}
\newcommand{\projektseminar}[1]{%
	\kopfeintrag{Projektseminar}{#1}{#1}
}
\newcommand{\oberseminar}[1]{%
	\kopfeintrag{Oberseminar}{#1}{#1}
}
\newcommand{\veranstaltung}[1]{%
	\kopfeintrag{Veranstaltung}{#1}{#1}
}
\newcommand{\kopfeintrag}[3]{%
    {#1:}&{\Large\bfseries #2}%
    \fakesubsection{#3}\\[3mm]
}%#1=Veranstaltung; #2=lfd. Seite; #3: Inhalts-Verz.
%
\newcommand{\untertitel}[1]{%
    & {\large\bfseries (#1)}\\[3mm]
}
\newcommand{\leereintrag}[2]{%
    {#1:}&{\bfseries  #2}\\[3mm]
}
\newcommand{\dozent}[1]{%
    {Dozent:}&{\bfseries  #1}\\[3mm]
}
\newcommand{\dozentin}[1]{%
    {Dozentin:}&{\bfseries  #1}\\[3mm]
}
\newcommand{\dozenten}[1]{%
    {Dozenten:}&{\bfseries  #1}\\[3mm]
}
\newcommand{\zeitort}[1]{%
    {Zeit/Ort:}& {\bfseries  #1}\\[3mm]
}
\newcommand{\uebungen}[1]{%
    {{\"U}bungen:}& {\bfseries  #1}\\[3mm]
}
\newcommand{\tutorium}[1]{%
    {Tutorium:}& {\bfseries  #1}\\[3mm]
}
\newcommand{\vorbesprechung}[1]{%
    {Vorbesprechung:}& {\bfseries  #1}\\[3mm]
}
\newcommand{\teilnehmerliste}[1]{%
    {Teilnehmerliste:}&{#1}\\[3mm]
}
\newcommand{\fragestunde}[1]{%
    {Fragestunde:}&{#1} \\[3mm]
}
\newcommand{\webseite}[1]{%
    {Web-Seite:}&{#1} \\[3mm]
}
\newcommand{\artderveranstaltung}[1]{%
    {Art der Veranstaltung:}&{#1}\\
}
\newcommand{\typischessemester}[1]{\relax}
%    {Typisches Semester:}&{#1}\\} % gestrichen 02.08.2016

\newcommand{\ectspunkte}[1]{%
    {ECTS-Punkte:}&{#1 Punkte}\\
}
%
\newcommand{\mastermodulrein}[1]{\relax}
% %ab WS2011/12 {Master-Studiengang:}&{geeignet f\"ur das Modul \emph{Reine Mathematik}}\\}
\newcommand{\mastermodulangew}[1]{\relax}
% % %ab WS2011/12 {Master-Studiengang:}&{geeignet f\"ur das Modul \emph{Angewandte Mathematik}}\\}
\newcommand{\mastermodulreinangew}[1]{\relax}
% %ab WS2011/12 {Master-Studiengang:}&{geeignet f\"ur die Module \emph{Reine Mathematik} und \emph{Angewandte Mathematik}}\\ }

\newcommand{\studienschwerpunkt}[1]{\relax}
% % % \newcommand{\studienschwerpunkt}[1]{%
% % %     {Studienschwerpunkt:}&{#1}\\
% % % }
\newcommand{\notwendigevorkenntnisse}[1]{%
    {Notwendige Vorkenntnisse:}&{#1}\\
}
\newcommand{\nuetzlichevorkenntnisse}[1]{%
    {N{\"u}tzliche Vorkenntnisse:}&{#1}\\
}
\newcommand{\pruefungsrelevan}[1]{\relax}%
%    {Pr{\"u}fungsrelevanz:}&{#1}\\} % gestrichen 02.08.2016

\newcommand{\folgeveranstaltungen}[1]{%
    {Folgeveranstaltungen:}&{#1}\\
}
\newcommand{\zusatzkommentar}[1]{\relax}%
%   {Kommentar:}&{{#1}}\\} gestrichen 02.08.2016
%
\newcommand{\zusatzbemerkung}[1]{%
   {Bemerkung:}&{{#1}}\\
}
\newcommand{\studienpruefungsleistung}{Studien-/Prüfungsleistung: & Die Anforderungen an Studien- und Prüfungsleistungen entnehmen Sie bitte dem aktuellen Modulhandbuch Ihres Studiengangs.\\}
%
\newcommand{\studienleistung}[1]{\relax}%
% {Studienleistung:}&{{#1}}\\} % gestrichen 02.08.2016
%
\newcommand{\pruefungsleistung}[1]{\relax}%
% {Pr{\"u}fungsleistung:}&{{#1}}\\} % gestrichen 02.08.2016
   
\newcommand{\sprechstundedozent}[1]{\relax}%
%     {Sprechstunde Dozent:}&{#1}\\} % gestrichen 02.08.2016
%
\newcommand{\sprechstundedozentin}[1]{\relax}%
%     {Sprechstunde Dozentin:}&{#1}\\} % gestrichen 02.08.2016
%
\newcommand{\sprechstundedozenten}[1]{\relax}%
%     {Sprechstunde Dozenten:}&{#1}\\} % gestrichen 02.08.2016
%
\newcommand{\sprechstundeassistent}[1]{\relax}%
%     {Sprechstunde Assistent:}&{#1}\\} % gestrichen 02.08.2016
%
\newcommand{\sprechstundeassistentin}[1]{\relax}%
%     {Sprechstunde Assistentin:}&{#1}\\} % gestrichen 02.08.2016
%
\newcommand{\zusaetzlicherkommentar}[1]{\relax}%
%    {Kommentar:}&{#1}\\} % gestrichen 02.08.2016
%
\newcommand{\zusaetzlichebemerkung}[1]{%
    {Bemerkung:}&{#1}\\
}
\newcommand{\verwendbarkeit}[1]{%
    {Verwendbarkeit:}&{#1}\\
}
\newcommand{\inhalt}[1]{{%
        {\large\textbf{Inhalt:}}\\[.5ex]%
         #1}
}
\newcommand{\fakesection}[1]{%
	\phantomsection%
	\addcontentsline{toc}{section}{#1}%
}
\newcommand{\fakesubsection}[1]{\phantomsection%
	\addcontentsline{toc}{subsection}{#1}%
}





% % \newcommand{\Beff}[2]{\verb+#1+}{{#1}{#2}}
% % % \newcommand{\vorlesung}[1]{%
% % % 	\kopfeintrag{Vorlesung}{#1}{#1}
% % \newcommand{\Befff}[3]{\verb+#1+}{{#2}{#3}}
% % % \newcommand{\leereintrag}[2]{%
% % %     {#1:}&{\bfseries  #2}\\[3mm]
% % \newcommand{\Beffff}[4]{\verb+#1+}{{#2}{#3}{#4}}
% % % \newcommand{\kopfeintrag}[3]{%
% % %     {#1:}&{\Large\bfseries #2}%
%     \fakesection{#3}\\[3mm]
